%------------------------------------------------------------------------------%
\begin{exercises}

%------------------------------------------------------------------------------%
\begin{exercise}
Nos itens a seguir, esboce a região delimitada pelas curvas indicadas
e calcule a sua área.
\begin{mathlist}
  \item y=e^x, \quad y=x^2-1, \quad x=-1, \quad x=1
  \item y=\sin(x), \quad y=x, \quad x=\frac{\pi}{2}, \quad x=\pi
  \item y=x^2-2x, \quad y=x+4
  \item y=\frac{1}{x}, \quad y=\frac{1}{x^2}, \quad x=2
  \item x=1-y^2, \quad x=y^2-1
  \item 4x+y^2=12, \quad x=y
  \item y=\tan(x), \quad y=2\sin(x), \quad -\frac{\pi}{3} \leq x \leq \frac{\pi}{3}
  \item y=x^3, \quad y=x
  \item y= \abs{x}, \quad y=x^2-2
  \item y=\frac{1}{x}, \quad y=x, \quad y=\frac{1}{4}x, \quad x>0
  \item y=\cos(x), \quad y=\sin(2x), \quad x=0, \quad x=\frac{\pi}{2}
\end{mathlist}
\begin{answer}
\begin{mathlist}
  \item e-\frac{1}{e}+ \frac{4}{3}
  \item \frac{3\pi^2}{8}-1
  \item \frac{125}{3}
  \item \ln 2-\frac{1}{2}
  \item \frac{8}{3}
  \item \frac{64}{3}
  \item 2-2\ln 2
  \item \frac{1}{2}
  \item \frac{20}{3}
  \item \ln 2
  \item \frac{1}{2}
\end{mathlist}
\end{answer}
\end{exercise}

%------------------------------------------------------------------------------%
\begin{exercise}
Use o Cálculo para encontrar a área do triângulo com vértices
$(0, 0), (3, 1)$ e $(1, 2)$.
\begin{answer}
  $A= \frac{5}{2}$
\end{answer}
\end{exercise}

%------------------------------------------------------------------------------%
\begin{exercise}
Calcule a área entre as curvas $y=x^2e^{-x}$ e $y=xe^{-x}$.
\begin{answer}
  $A=3e^{-1}-1$
\end{answer}
\end{exercise}

%------------------------------------------------------------------------------%
%\begin{exercise}
%  A parábola $y= \frac{1}{2}x^2$ divide o disco $x^2+y^2 \leq 8$ em duas partes.
%  Encontre a área dessas duas partes.
%\end{exercise}

%------------------------------------------------------------------------------%
\begin{exercise}
  Um arquiteto propôs colocar no piso de uma praça pastilhas coloridas
  em uma região plana cujo o formato se assemelha (respeitando a escala)
  a região delimitada simutâneamente pelas quarto curvas
  $y= ln(x)$, $y=e^x$, $y=1-x$ e $x=e$.
  Esboce a região e calcule a área.
\end{exercise}

%------------------------------------------------------------------------------%
\begin{exercise}
Determine a área da região limitada exatamente pelas três curvas
$y =x^2$, $ y = 2- x^2$ e $y = 2x + 8$.
\begin{answer}
  $\frac{100}{3}$
\end{answer}
\end{exercise}

%------------------------------------------------------------------------------%
\begin{exercise}
  Encontre os valores de $c$ tais que a área da região delimitada pelas
  parábolas $y = x^2 - c^2$ e $y=c^2-x^2$ seja $\frac{8}{3}$.
\begin{answer}
  $c=1$ e $c=-1$
\end{answer}
\end{exercise}

%------------------------------------------------------------------------------%
\begin{exercise}
  A parábola $y=\frac{1}{2}x^2$ divide o disco $x^2+y^2 \leq 8$ em duas partes.
  Encontre a área de ambas as partes.
\end{exercise}

%------------------------------------------------------------------------------%
%\begin{exercise}
%\begin{answer}
%\end{answer}
%\end{exercise}

%------------------------------------------------------------------------------%
%\begin{exercise}
%\begin{mathlist}[2]
%  \item
%  \item
%  \item
%\end{mathlist}
%\begin{answer}
%  \begin{alphalist}[3]
%    \item
%    \item
%    \item
%  \end{alphalist}
%\end{answer}
%\end{exercise}

\end{exercises}
%------------------------------------------------------------------------------%

